
\svnid{$Id: $}
\svnid{$URL: $}

\section{Constant, time-varying, and interdependent (Q-h) boundary conditions.}
\newrefsegment

%---------------------------------------------------------------------------------
\section*{Quality Assurance}
\begin{tabular}{@{}p{27mm}@{}p{0mm}p{\textwidth-52mm-48pt}p{15mm}p{10mm}}
\textbf{Date} & \textbf{:} &   August 02 2019\\ 
\textbf{Author} & \textbf{:} & Christianne Luger  & \textbf{Initials:} & CL \\
\textbf{Review} & \textbf{:} & Bert Jagers & \textbf{Initials:} & BJ \ldots 
\end{tabular}

%---------------------------------------------------------------------------------
\section*{Version information}
{{
\begin{tabular}{@{}p{27mm}@{}p{0mm}p{\textwidth-27mm-24pt}}
\textbf{Executable}   & \textbf{:} & Deltares, D-Flow FM Version 1.2.59.64502M, Aug 02 2019, 10:44:19  \\
\textbf{Location input}     & \textbf{:} & \\
\textbf{Revision input} & \textbf{:} &  \\
\textbf{Location output}     & \textbf{:} & \url{https://repos.deltares.nl/repos/DSCTestbench/trunk/cases/e02_dflowfm/f101_1D-boundaries/c01_steady-state-flow} \\
\textbf{Revision output} & \textbf{:} & 
\end{tabular}
}}

%---------------------------------------------------------------------------------
\section*{Purpose}
The purpose of this validation test is to verify the proper functioning of constant and time-varying boundary conditions for upstream discharge and downstream water level in 1D \dflowfmfull simulations.
Also the "Q-h relationship" to force the water level based on simulated discharge is tested.

%---------------------------------------------------------------------------------
\section*{Linked claims}
Claims that are investigated in the current test case are:
\begin{itemize}
\item \dflowfm can accurately simulate steady and unsteady hydrostatic flow on 1D networks.
\end{itemize}

%---------------------------------------------------------------------------------
\section*{Approach}
Four straight channels, T1 to T4, (see \autoref{fig:10101:Test model lay-out}) are parametrized with different boundary condition types, namely constant in time, time-varying, and interdependent (Q-h).
The latter is specified using a direct relationship between discharge and the local water level.

The boundary conditions of all sub-models is such that they ultimately should converge to the same values.
The proper functioning of all boundary conditions types is confirmed if all sub-models reach the same water level and discharge end result and show the correct dynamic behaviour.

\begin{figure}[H]
    \centering
    \includegraphics*[width=0.95\textwidth]{figures/e02_f101_c01_layout.JPG}
    \caption{Lay-out of the test model}
		\label{fig:10101:Test model lay-out}
\end{figure}


%---------------------------------------------------------------------------------
\section*{Conclusion} 
Independent of the boundary condition type, all models reach the same equilibrium conditions.
Furthermore, the discharges recorded at the midpoint of the channels depict the correct and expected dynamic behaviour with respect to the upstream and downstream boundary conditions parametrized.
Therefore, it is verified through this test that constant, time-varying and interdependent (Q-h) boundary conditions function properly in 1D \dflowfmfull simulations.

%---------------------------------------------------------------------------------
\section*{Model description}
The test comprises of four separate sub-models T1 to T4 (see \autoref{fig:10101:Test model lay-out}). 
The models are 2 km in length and have an equidistant grid spacing of dx=50m. 
Roughness is constant at a Ch\'ezy roughness coefficient of 45 m\textsuperscript{0.5}s\textsuperscript{-1} throughout the length and across the cross section of the sub-models as the lumped conveyance approach is used (see \autoref{10101:description of Lumped conveyance approach}). 
The same YZ cross sectional profile is applied uniformly throughout all sub-models (see \autoref{fig:10101:YZCrossSection}), but the bed levels decrease by 0.5 m from chainage=0 (east) to chainage=2000 (west) for a slope of 0.0025. 

\begin{Note}
The kernel of \dflowfm does not allow vertical line segments in the YZ-profile. 
Hence, vertical line segments in this trapezoidal channel at 0 and 50 m are tilted slightly by 0.001 m.
\end{Note}

\begin{figure}[H]
    \centering
    \includegraphics*[width=0.95\textwidth]{figures/e02_f101_c01_cross_section.jpg}
    \caption{YZ Cross-section definition}
    \label{fig:10101:YZCrossSection}
\end{figure}

The four sub-models are parametrized with different types of upstream and downstream boundary condition which ultimately converge to the same values (see \autoref{10101:Boundary conditions}). 

\begin{longtable}{|p{15mm}|p{20mm}|p{27mm}|p{70mm}|}
\caption{Boundary conditions applied in sub-models T1 to T4.}
\label{10101:Boundary conditions} \\ 
\hline
\STRUT Sub-model & Chainage \unitbrackets{m} & ID & Boundary condition \\ \hline
\endhead
\hline
\endfoot
%
\STRUT T1 & 0.00    & T1\_Up\_Bnd  & Constant discharge of 100 m\textsuperscript{3}s\textsuperscript{-1}\\
\STRUT    & 2000.00 & T1\_Dwn\_Bnd & Constant water level of 2.5 m\\ \hline
\STRUT T2 & 0.00    & T2\_Up\_Bnd  & Time-varying Discharge (See \autoref{fig:10101:Qt})\\
\STRUT    & 2000.00 & T2\_Dwn\_Bnd & Constant water level of 2.5 m\\ \hline
\STRUT T3 & 0.00    & T3\_Up\_Bnd  & Constant discharge of 100 m\textsuperscript{3}s\textsuperscript{-1}\\
\STRUT    & 2000.00 & T3\_Dwn\_Bnd & Time-varying Water Level (See \autoref{fig:10101:Ht})\\ \hline
\STRUT T4 & 0.00    & T4\_Up\_Bnd  & Time-varying Discharge (See \autoref{fig:10101:Qt})\\
\STRUT    & 2000.00 & T4\_Dwn\_Bnd & Q-h boundary condition (See \autoref{10101:Qhbnd})\\
\end{longtable}

\begin{figure}[H]
    \centering
    \includegraphics*[width=0.95\textwidth]{figures/e02_f101_c01_Qt.jpg}
    \caption{Time-Varying Discharge Boundary Condition for T2 and T4}
    \label{fig:10101:Qt}
\end{figure}

\begin{figure}[H]
    \centering
    \includegraphics*[width=0.95\textwidth]{figures/e02_f101_c01_Ht.jpg}
    \caption{Time-Varying Water Level Boundary Condition for T3}
    \label{fig:10101:Ht}
\end{figure}

\begin{longtable}{|p{30mm}|p{40mm}|}
\caption{Discharge-Water Level Interdependent Boundary Condition for T4}
\label{10101:Qhbnd} \\
\hline
\STRUT Discharge [m\textsuperscript{3}s\textsuperscript{-1}] & Water Level [m] \\ \hline \endhead
\hline
\endfoot
%
\STRUT 50.00 & 1.25 \\ \hline
\STRUT 100.00 & 2.50 \\ \hline
\STRUT 150.00 & 3.75 \\
\end{longtable}


%---------------------------------------------------------------------------------
\section*{Results}

The modeled water depths and discharges along sub-models \textbf{T1} and \textbf{T4} at the end of the simulation (72 hours) are on \autoref{fig:10101:WaterDepthEntireLength} and \autoref{fig:10101:DischargeEntireLength} respectively.
\autoref{fig:10101:DischargeXsection} depicts the discharges throughout the simulation period recorded at observation cross-sections at the middle of the sub-models (see \autoref{fig:10101:Test model lay-out}).

\begin{figure}[H]
    \centering
    \includegraphics*[width=0.95\textwidth]{figures/e02_f101_c01_Water_Depth_EntireLength.jpg}
    \caption{Modeled water depths at the last time step (72 hours)}
    \label{fig:10101:WaterDepthEntireLength}
\end{figure} 

\begin{figure}[H]
    \centering
    \includegraphics*[width=0.95\textwidth]{figures/e02_f101_c01_Discharge_EntireLength.jpg}
    \caption{Modeled discharge at the last time step (72 hours)}
    \label{fig:10101:DischargeEntireLength}
\end{figure} 

\begin{figure}[H]
    \centering
    \includegraphics*[width=0.95\textwidth]{figures/e02_f101_c01_Discharge_ObsCross.jpg}
    \caption{Modeled discharge for all time-steps at observation cross-sections}
    \label{fig:10101:DischargeXsection}
\end{figure} 


%---------------------------------------------------------------------------------
\section*{Analysis of results}
All sub-models converge effectively to the same end results for discharges and water levels.
The maximum differences are negligible.

All sub-models further exhibit the expected dynamic behaviour. 
The discharges reach the same value after 30 hours which is the same time that the boundary conditions of the sub-models are parametrized to converge.
Sub-models T1 and T3 are forced with a constant upstream discharge of 100 m\textsuperscript{3}s\textsuperscript{-1}.
The observed discharge at the midpoint matches the specified value throughout the simulation apart from a brief ramp-up period from the dry starting conditions.
In sub-model T3 the discharge deviates slightly from the specified discharge value until 30 hours as the sub-model adjusts continuously to the downstream time-varying water level boundary conditions -- as can be observed from the final water depth curve (see \autoref{fig:10101:WaterDepthEntireLength}) the backwater effect of the downstream boundary reaches up to the upstream boundary.
Sub-model T2 and T4 are both forced using a time-varying upstream discharge boundary condition that ramps up to 100 m\textsuperscript{3}s\textsuperscript{-1} after 30 hours.
This time-variation is reflected by the discharge observed halfway the channels.
The initial negative discharge values at the beginning of the simulation for these sub-models are caused by water flowing upstream from the downstream water level boundary.

%---------------------------------------------------------------------------------
\subsection{Lumped conveyance approach} \label{10101:description of Lumped conveyance approach}
In the lumped conveyance approach, it is assumed that there is an uniform flow velocity (\=U) in the cross-sectional profile for a given water level (see \autoref{fig:10101:Lumped approach}).
Hence, it is assumed that differences in roughness coefficients and flow velocities across the cross-sectional profile can be neglected.

\begin{figure}[H]
    \centering
    \includegraphics*[width=0.65\textwidth]{figures/e02_f101_c01_lumped_conveyance.png}
    \caption{Concept of the lumped conveyance approach}
    \label{fig:10101:Lumped approach}
\end{figure}

The lumped conveyance $K$ is computed as a function of water depth $h$:

\begin{equation}
K(h) = A(h)C(h)\sqrt{R(h)}
\end{equation}

where:
\begin{itemize}
\item $K(h)$ Lumped conveyance at $h$
\item $A(h)$ Cross-sectional area at $h$
\item $C(h)$ Ch\'ezy friction value at $h$
\item $P(h)$ Wetted perimeter at $h$
\item $R(h)$ Hydraulic radius at $h(R(h)=A(h)/P(h))$
\end{itemize}

Note that the wetted perimeter $P$ as function of $h$ would be discontinuous when a part of the profile were horizontal; therefore, we don't allow for that situation.
Roughness types other than Ch\'ezy will be converted to their Ch\'ezy equivalent based on the local water depth $h$ and hydraulic radius $R(h)$.

\printrefsegment

